% prologue.tex --- Prologue (Scenes 1--3)
% Inputted by lecture.tex; assumes the AINITAcademic theme is loaded
% and \title/\subtitle/\author/\institute/\date are already set.

% ---------------------------------------------------------------------------
% Scene 1 --- Title slide
% ---------------------------------------------------------------------------
\renewcommand{\AINITsection}{}
\AINITmaketitle

% ---------------------------------------------------------------------------
% Scene 2 --- Funes opener
% Full-page italic epigraph; attribution in small caps, mutedWarmGray.
% Empty third argument suppresses the section title underneath.
% ---------------------------------------------------------------------------
\sectiondividerwithportrait%
  {borges.jpg}%
  {``To think is to forget differences, generalize, abstract.''}%
  {Jorge Luis Borges, \textit{Funes the Memorious}, 1942}%
  {}%

% ---------------------------------------------------------------------------
% Scene 2b --- Bridge from Week 1: the transformer's two halves
% Recall the encoder/decoder summary from w1_2_llm.pdf, slide 4 (Ge'ez games).
% Left column: short prose. Right column: minimal TikZ block diagram.
% ---------------------------------------------------------------------------
\begin{frame}{Last week, in language}
  \begin{columns}[T]
    \column{0.62\textwidth}
    \vspace{0.3em}
    You met the transformer. The summary we landed on:

    \medskip

    \emph{An LLM \(\approx\) encoder (builds meaning) \(+\) decoder
    (continues the sentence).}

    \medskip

    The \AINITemph{encoder} reads a prefix and produces an
    \AINITemph{embedding} \textemdash{} a vector summarising what the
    prefix says. The \AINITemph{decoder} takes that vector and writes a
    probability distribution over the next token.

    \column{0.34\textwidth}
    \vspace{0.5em}
    \centering
    \begin{tikzpicture}[
      node distance=0.3cm,
      data/.style={
        font=\itshape\footnotesize, color=mutedWarmGray,
        align=center, inner sep=2pt
      },
      block/.style={
        rectangle, rounded corners=1.5pt,
        draw=accentOxblood, line width=0.5pt,
        minimum height=0.7cm, minimum width=1.8cm,
        fill=bgCream, align=center,
        font=\itshape\small
      },
      ar/.style={->, line width=0.5pt, color=accentOxblood, >=Latex}
    ]
      \node[data] (prefix) {prefix};
      \node[block, below=of prefix] (enc) {encoder};
      \node[data, below=of enc] (vec) {embedding\\(vector)};
      \node[block, below=of vec] (dec) {decoder};
      \node[data, below=of dec] (probs) {probabilities\\over next token};
      \draw[ar] (prefix) -- (enc);
      \draw[ar] (enc) -- (vec);
      \draw[ar] (vec) -- (dec);
      \draw[ar] (dec) -- (probs);
    \end{tikzpicture}
  \end{columns}
\end{frame}

% ---------------------------------------------------------------------------
% Scene 2c --- Reframe: encoder = compression, decoder = prediction
% ---------------------------------------------------------------------------
\begin{frame}{Two halves of one machine}
  The encoder and decoder have older names.

  \medskip
  \begin{itemize}
    \setlength{\itemsep}{0.5em}
    \item Encoder \(\approx\) \AINITemph{compression}.
          A long sequence becomes a short vector.
    \item Decoder \(\approx\) \AINITemph{prediction}.
          The vector becomes a bet about what comes next.
  \end{itemize}

  \medskip
  For ninety minutes today we will do each half explicitly \textemdash{}
  but for data that isn't language. \AINITemph{Customers, not words.
  Campaigns, not tokens.}

  \medskip
  This framework is older than the transformer. It is from a 1948 paper
  by Shannon.
\end{frame}

% ---------------------------------------------------------------------------
% Scene 3 --- Thesis paragraph
% Standard content slide, left-aligned, no bullets.
% ---------------------------------------------------------------------------
\begin{frame}{The promise}
  Machine learning is the engineering of two old theorems.
  \AINITemph{Compression} \textemdash{} the shortest faithful description of
  the past \textemdash{} and \AINITemph{prediction} \textemdash{} the best
  bet about the future \textemdash{} are two views of the same subject,
  governed by Shannon's 1948 theorems. We will build them in two halves and
  collapse them in the third.
\end{frame}
